% Vietnamese translation for OI Public Library
% Source-File: IOI中国国家候选队论文/2006/贾由/贾由.ppt
\documentclass[11pt]{article}
\usepackage{oipl-vn}

\title{Tối ưu thuật toán qua bài toán đồ thị\\{\large Ghi chú theo slide}}
\author{Jia You}
\date{2006}

\begin{document}
\maketitle

\origtitle{Algorithm optimization through graph problems}
\sourcepath{IOI China candidate team papers / 2006 / Jia You / Jia You.ppt}

\section{Phạm vi}

Bộ slide đi kèm bài viết về tối ưu thuật toán từ góc nhìn các bài toán đồ
thị. Ghi chú này dựa trên bản bài viết đã dịch.

\section{Tìm điểm đặc biệt}

Một thuật toán đồ thị kinh điển thường giải mô hình tổng quát. Đề thi lại hay
thêm cấu trúc đặc biệt: hai phía, cây, bậc nhỏ, trọng số đơn vị, thứ tự topo,
hoặc ràng buộc hình học. Tối ưu tốt bắt đầu bằng việc phát hiện cấu trúc ấy và
loại các phần của thuật toán tổng quát không còn cần thiết.

\section{Ví dụ ghép hai phía}

Ghép cực đại hai phía có thể quy về luồng cực đại, nhưng mạng thu được có
dạng rất riêng: mọi sức chứa là một, đường tăng bắt đầu ở một đỉnh chưa ghép
bên trái và kết thúc ở một đỉnh chưa ghép bên phải. Khi bỏ nguồn, đích và các
cạnh hiển nhiên, ta đi tới tư tưởng của thuật toán Hungary.

\section{Sửa thuật toán sai}

Bài viết cũng nhấn mạnh một con đường tối ưu khác: bắt đầu từ thuật toán chưa
đúng hoặc chưa đủ nhanh, phân tích phản ví dụ, rồi thêm điều kiện còn thiếu.
Quá trình này thường giúp nhận ra chính xác phần thông tin mà trạng thái hoặc
cấu trúc dữ liệu cần giữ.

\section{Ghi nhớ}

Tối ưu hóa không chỉ là thay hàng đợi bằng heap hay dùng bitset. Với bài đồ
thị, câu hỏi quan trọng hơn là: đồ thị của đề khác đồ thị tổng quát ở đâu, và
sự khác biệt đó cho phép bỏ bớt những thao tác nào?

\end{document}
